% ============================================================
%  ACCESSIBILITY RELAXATION AND FRIEDMANN INSTABILITY:
%  AN RSVP REINTERPRETATION OF COSMOLOGICAL SADDLE GEOMETRY
%  Flyxion — Independent Researcher — May 2026
% ============================================================

\documentclass[12pt,a4paper]{article}

\usepackage[T1]{fontenc}
\usepackage{palatino}

\usepackage{geometry}
\geometry{top=1.15in, bottom=1.15in, left=1.25in, right=1.25in}

\usepackage{microtype}
\usepackage{setspace}
\onehalfspacing

\usepackage{amsmath,amssymb,amsthm,mathtools}
\usepackage{bm}
\usepackage{enumitem}
\usepackage{graphicx}
\usepackage{booktabs}
\usepackage{natbib}
\bibpunct{[}{]}{,}{n}{}{;}

\usepackage{titlesec}
\titleformat{\section}{\normalfont\large\bfseries}{\thesection.}{0.6em}{}
\titleformat{\subsection}{\normalfont\normalsize\bfseries}{\thesubsection.}{0.5em}{}
\titleformat{\subsubsection}{\normalfont\normalsize\itshape}{\thesubsubsection.}{0.5em}{}

\usepackage[colorlinks=true,linkcolor=black,citecolor=black,urlcolor=black]{hyperref}

% ---- Theorem environments ----
\theoremstyle{plain}
\newtheorem{theorem}{Theorem}[section]
\newtheorem{proposition}[theorem]{Proposition}
\newtheorem{lemma}[theorem]{Lemma}
\newtheorem{corollary}[theorem]{Corollary}

\theoremstyle{definition}
\newtheorem{definition}[theorem]{Definition}
\newtheorem{principle}[theorem]{Reduction Principle}
\newtheorem{assumption}[theorem]{Assumption}

\theoremstyle{remark}
\newtheorem{remark}[theorem]{Remark}

% ---- Notation macros ----
\newcommand{\Phic}{\Phi_c}
\newcommand{\Tacc}{\mathcal{T}}
\newcommand{\Ggroup}{\mathcal{G}}
\newcommand{\Uacc}{\mathcal{U}}
\newcommand{\Sacc}{\mathcal{S}}
\newcommand{\piproj}{\pi}
\newcommand{\Hacc}{\mathcal{H}}
\newcommand{\Lacc}{\mathcal{L}}
\newcommand{\nabacc}{\nabla^{\!\mathrm{acc}}}
\newcommand{\Delacc}{\Delta_{\mathrm{acc}}}
\newcommand{\SM}{\mathit{SM}}
\newcommand{\Mink}{\mathit{M}}
\newcommand{\Ffam}{\mathcal{F}}
\newcommand{\dclos}{d_{\mathrm{closure}}}
\newcommand{\etaeff}{\eta_{\mathrm{eff}}}
\newcommand{\Cbz}{\mathcal{C}_{\beta z}}
\newcommand{\bcurl}{\beta_{\mathrm{curl}}}
\newcommand{\lAn}[1]{\lambda^A_{#1}}
\newcommand{\lBn}[1]{\lambda^B_{#1}}
\newcommand{\lA}{\lambda^A}
\newcommand{\lB}{\lambda^B}
\newcommand{\kA}{\kappa_A}

% ---- Title ----
\title{%
  \LARGE\bfseries
  Accessibility Relaxation and Friedmann Instability\\[6pt]
  \large
  An RSVP Reinterpretation of Cosmological Saddle Geometry
}
\author{
  Flyxion\\[4pt]
  \normalsize Independent Researcher
}
\date{May 2026}

\begin{document}

\maketitle

\begin{abstract}
We develop a reinterpretation of cosmological dynamics within the Relativistic
Scalar--Vector Plenum (RSVP) framework, establishing a precise mathematical
relationship between the instability structure of Friedmann spacetimes and the
accessibility relaxation geometry of the RSVP lamphrodyne system. The central
result of Alexander, Temple, and Vogler --- that the critical Friedmann spacetime
is an unstable saddle rest point whose infinite-order instability hierarchy is
controlled by a second-order eigenvalue condition --- is recast as a theorem about
a constrained accessibility bundle in RSVP trajectory space. We introduce the
\emph{closure depth} $\dclos$ as a new invariant of accessibility bundles and
prove that the Friedmann instability satisfies $\dclos(\Phic) = 2$, making it a
depth-two object in the CLIO constraint-closure hierarchy. The RSVP lamphrodyne
functional, derived from an accessibility metric and a divergence constraint,
reduces to the Alexander--Temple--Vogler STV-PDE under five explicit conditions
constituting a Reduction Principle. Relaxing these conditions generates three
families of physical corrections. The third-order redshift--luminosity coefficient
is predicted to satisfy $C_{\mathrm{RSVP}} = 0.3591 + \etaeff + 3\mu^2\alpha_\Phi^2/4
+ \bcurl\cdot\mathcal{F} + O(\etaeff^2,\mu^3,\bcurl^2)$, exceeding both the
Alexander--Temple--Vogler value of $+0.3591$ and the $\Lambda$CDM value of
$-0.1804$. When the curl coupling $\bcurl > 0$, the same entropy--vorticity source
$\nabla S\cdot\omega$ that perturbs the redshift coefficient also sources parity-odd
birefringence, predicting a nonzero redshift--birefringence cross-correlation
$\Cbz \neq 0$ absent from both $\Lambda$CDM and the Alexander--Temple--Vogler
system.
\end{abstract}

\tableofcontents
\newpage

% ============================================================
\section{Introduction}
\label{sec:intro}
% ============================================================

The discovery that the critical Friedmann spacetime is an unstable saddle point
within the solution space of Einstein's original field equations, without any
cosmological constant, constitutes one of the more consequential recent results
in mathematical cosmology. The programme establishing this instability developed
across three papers~\cite{Smoller2012,STV2017,ATV2026}: Smoller and Temple~\cite{Smoller2012}
established the general relativistic self-similar wave framework; Smoller, Temple,
and Vogler~\cite{STV2017} introduced the STV-PDE and STV-ODE systems and identified
the instability structure; Alexander, Temple, and Vogler~\cite{ATV2026} provided
the definitive eigenvalue analysis proving that every Friedmann spacetime is unstable
to radial perturbation at every perturbative order, and that generic underdense
perturbations produce accelerated expansion as a natural consequence of this
instability rather than as a signature of dark energy. Their framework --- which we
call the ATV system --- provides a mathematically rigorous alternative to $\Lambda$CDM
for the anomalous acceleration problem, while remaining entirely within classical
general relativity.

The present paper does not dispute these results. It asks a different question:
what ontological status should be assigned to the objects in the ATV analysis ---
the saddle point, the unstable manifold, the gauge-fixing, the recursive eigenvalue
structure --- and is there a broader framework within which they appear as a
projected sector of a more general dynamical theory?

The present paper develops such a framework, which we call the
\emph{accessibility relaxation formalism}. The framework treats cosmological dynamics
as the macroscopic behaviour of a continuous medium described by three coupled fields:
a scalar accessibility potential $\Phi$, a lamphrodyne flow velocity $\mathbf{V}$,
and a configurational entropy density $S$. Gravitational binding, cosmic redshift,
and structure formation emerge from the interactions of these fields through a
variational principle --- the \emph{lamphrodyne functional} --- whose Lagrangian
is constrained by locality, rotational invariance, ghost-freedom, and a divergence
condition linking the longitudinal flow mode to scalar and entropy sources.
The framework connects to Jacobson's thermodynamic derivation of the Einstein
equation~\cite{Jacobson1995}, Verlinde's entropic gravity~\cite{Verlinde2011},
and Padmanabhan's thermodynamic cosmology~\cite{Padmanabhan2010} as special
limits, while extending them to a dynamical, non-equilibrium medium with explicit
attractor geometry.

The relationship between the accessibility relaxation formalism and the ATV system
is not merely interpretive. The ATV saddle point $\SM$, its eigenvalue spectrum,
the gauge decomposition effected by the ``time since the Big Bang'' fixing, and
the recursive closure of the unstable manifold all find precise counterparts in
the accessibility formalism introduced here. Specifically, we prove:

\begin{itemize}[leftmargin=1.5em]
\item The ATV gauge-fixing corresponds to choosing a section of the fiber bundle
  $\Tacc \to \Tacc/\Ggroup$, where $\Ggroup$ acts by time translation and the
  fiber direction is the eigenspace of the leading-order negative eigenvalue
  $\lBn{1} = -1$.

\item The ATV eigenvalue formula $\lAn{n} = 2n/3$, $\lBn{n} = (2n-5)/3$ is
  reproduced as the spectrum of the $n$-th level projected Hessian of the
  lamphrodyne functional, under five explicit reduction conditions.

\item The ATV theorem that infinite-order instability is controlled by a second-order
  condition is equivalent to the statement that the closure depth of the
  unstable accessibility bundle is $\dclos(\Phic) = 2$, a new invariant we
  introduce and compute.

\item Relaxing the reduction conditions generates corrections to the
  redshift--luminosity coefficient $C$ and sources parity-odd birefringence through
  a common entropy--vorticity coupling, producing a jointly falsifiable prediction
  absent from both ATV and $\Lambda$CDM.
\end{itemize}

The strategic value of this reinterpretation is twofold. First, it gives RSVP its
most direct contact with established relativistic mathematics: the ATV analysis
provides a rigorous local model inside a broader RSVP ontology, rather than
requiring RSVP to carry the observational burden of cosmology unaided. Second, it
extends the ATV framework by generating concrete corrections that are falsifiable
at the level of fourth-order supernova data and CMB polarization cross-correlations.

The paper is organised as follows. Section~\ref{sec:geometry} develops the RSVP
accessibility geometry and defines the lamphrodyne functional. Section~\ref{sec:gauge}
establishes the gauge decomposition and fiber structure. Section~\ref{sec:friedmann}
embeds the Friedmann saddle as an accessibility-critical rest point.
Section~\ref{sec:closure} proves the recursive closure theorem and introduces
the closure depth invariant, with its CLIO interpretation. Section~\ref{sec:lamp_limit}
derives the Reduction Principle and the ATV limit. Section~\ref{sec:redshift}
computes the lamphrodyne corrections to the redshift coefficient.
Section~\ref{sec:bire} establishes the coupled redshift--birefringence prediction.
Section~\ref{sec:disc} collects the observational discriminants. Section~\ref{sec:beyond}
discusses the broader programme of RSVP cosmology beyond expansion.

Throughout, standard results from ATV are cited by their theorem numbers in
\cite{ATV2026}. All RSVP objects are defined independently of ATV; the
correspondence is established as a theorem rather than assumed.

% ============================================================
\section{RSVP Accessibility Geometry}
\label{sec:geometry}
% ============================================================

The RSVP framework begins with a continuous medium whose state space carries more
structure than a single metric. The fundamental objects are three coupled fields
whose interactions generate the full range of physical phenomena --- gravitational,
thermodynamic, and cognitive --- through a single variational principle.

\begin{definition}[RSVP field triplet]
\label{def:triplet}
Let $\Omega \subset \mathbb{R}^3$ be an open bounded domain with Lipschitz
boundary and let $\mathcal{M} = \Omega \times (0,\infty)$. The RSVP medium is
described by three fields:
\begin{align}
  \Phi &\in H^1(\Omega), \quad \text{(scalar accessibility potential)},
  \label{eq:field_phi} \\
  \mathbf{V} &\in [H^1(\Omega)]^3, \quad \text{(lamphrodyne flow velocity)},
  \label{eq:field_V} \\
  S &\in L^2(\Omega)\cap L^\infty(\Omega),\; S\geq 0, \quad
    \text{(configurational entropy density)}.
  \label{eq:field_S}
\end{align}
In the cosmological projection, $\Phi$ plays the role of a scalar plenum potential
analogous to vacuum capacity, $\mathbf{V}$ encodes negentropic bulk flow, and $S$
drives redshift and constraint relaxation. The connection to emergent gravitational
thermodynamics follows Jacobson~\cite{Jacobson1995} and Padmanabhan~\cite{Padmanabhan2010}:
entropy gradients generate effective forces, and the accessible trajectory volume
in the medium provides the statistical underpinning for gravitational binding.
\end{definition}

\begin{remark}[Effective-theory status of the RSVP framework]
\label{rem:eft_status}
The accessibility relaxation formalism should presently be understood as an
\emph{effective field theory} rather than a completed microscopic description
of spacetime. Its purpose is not to replace General Relativity at every scale
but to provide a dynamical framework within which the instability results of
Alexander, Temple, and Vogler can be embedded and extended. The scalar
accessibility field $\Phi$, lamphrodyne flow $\mathbf{V}$, and entropy density
$S$ are best regarded as macroscopic variables describing the relaxation state
of the medium under coarse-graining at the relevant cosmological scale.

In this reading, the RSVP Lagrangian is the most general effective action
consistent with the stated symmetry and stability axioms at two-derivative order.
Higher-derivative corrections, microscopic derivation of the coupling constants
$(\kappa, \alpha, \bcurl, \mu)$, and the relationship of the plenum fields to
any underlying quantum degrees of freedom all remain open problems. The programme
is in this respect analogous to the effective field theory of inflation or to
analogue-gravity models~\cite{MisnerThorneWheeler1973,HawkingEllis1973}: the
effective description is internally consistent and predictive within its domain,
without yet possessing a complete ultraviolet completion.
\end{remark}

\begin{definition}[Accessibility metric]
\label{def:acc_metric}
The \emph{accessibility metric} is the field-dependent deformation of the
background Riemannian metric $g_{ij}$ defined by
\begin{equation}
  g^{\mathrm{acc}}_{ij} = g_{ij} + \kappa\,\partial_i\Phi\,\partial_j\Phi,
  \qquad \kappa \geq 0.
  \label{eq:acc_metric}
\end{equation}
When $\kappa > 0$, the accessibility metric assigns greater length to
displacements along the gradient of $\Phi$, penalising motion in directions of
rapid scalar variation. This is the geometric realisation of the RSVP principle
that entropy gradients compress accessible trajectory space: regions of high
$|\nabla\Phi|$ are metrically distant in the accessibility sense. Operators
carrying the superscript $\mathrm{acc}$ --- gradient $\nabacc$, divergence
$\nabla^{\!\mathrm{acc}}\cdot$, curl $\nabacc\times$, and Laplacian $\Delacc$ ---
are computed with respect to $g^{\mathrm{acc}}$.
\end{definition}

\subsection{Effective-Theory Status of the Accessibility Metric}
\label{subsec:effective_metric}

The accessibility metric of Definition~\ref{def:acc_metric} should not be
interpreted as a fundamental modification of spacetime geometry. It is introduced
as an effective metric on accessibility space, analogous to the emergent metrics
appearing in acoustic gravity~\cite{MisnerThorneWheeler1973} and thermodynamic
state-space geometry. Its role is to encode how gradients of the scalar field
modify the volume of admissible trajectories available to the medium.

Equation~\eqref{eq:acc_metric} is the unique lowest-order rotationally invariant
rank-one deformation of $g_{ij}$ generated by the scalar field gradient at
two-derivative order. The present paper does not claim a microscopic derivation
of this structure. The framework should accordingly be understood as an effective
field theory whose metric structure is justified by symmetry counting and
phenomenological consistency, with derivation from a deeper microphysics of
accessibility remaining an open problem.

\begin{definition}[Lamphrodyne functional]
\label{def:lamp_functional}
The \emph{lamphrodyne functional} is
\begin{equation}
  \Lacc[\Phi,\mathbf{V}]
  = \int_{\Omega}
    \Bigl[
      \tfrac{1}{2}\,g^{\mathrm{acc},ij}\,\partial_i\Phi\,\partial_j\Phi
      + \tfrac{\alpha}{2}\,|\nabacc\cdot\mathbf{V}|^2
      + \tfrac{\bcurl}{2}\,|\nabacc\times\mathbf{V}|^2
      + \mu\,\mathbf{V}\cdot\nabacc\Phi
      + U(\Phi)
    \Bigr]
    d\mathrm{vol}_{g^{\mathrm{acc}}},
  \label{eq:lamp_functional}
\end{equation}
where $\alpha, \bcurl, \mu > 0$ are coupling constants, $U \in C^2(\mathbb{R})$
satisfies $U'' \geq 0$, and $g^{\mathrm{acc},ij}$ denotes the inverse of the
accessibility metric. The four terms encode, respectively: accessibility gradient
energy, compression resistance of the longitudinal flow mode, curl energy of the
transverse flow mode, and scalar--vector coupling. The potential $U(\Phi)$ provides
attractor geometry and may encode a MOND-like low-acceleration regime in the
appropriate limit, following the entropic force construction of Verlinde~\cite{Verlinde2011}.
\end{definition}

\subsection{Mode Stability and Derivative Order}
\label{subsec:ghosts}

The lamphrodyne functional~\eqref{eq:lamp_functional} contains first-derivative
couplings between $\Phi$ and $\mathbf{V}$ but no higher-order time derivatives
of any field. The classical Ostrogradsky instability associated with
non-degenerate higher-derivative Lagrangians therefore does not arise at the
level of the present theory: all kinetic terms are at most first order in $\partial_t$.

The coupling $\mu\mathbf{V}\cdot\nabacc\Phi$ is a first-order spatial coupling,
not a higher-time-derivative term, and its contribution to the Hessian is bounded
by standard Young's inequality estimates (as used in the proof of Lemma~\ref{lem:selfadjoint}).
A complete Hamiltonian decomposition of the scalar, longitudinal, and transverse
sectors into physical and unphysical modes --- analogous to the Coulomb-gauge
decomposition in constrained field theories~\cite{Arnold1989} --- is beyond the
scope of the present paper. The framework should therefore be regarded as an
effective field theory whose stability in the full nonlinear regime has not yet
been established, and a complete nonlinear stability analysis remains a priority
for future work.

\begin{definition}[Lamphrodyne flow equations]
\label{def:lamp_flow}
The \emph{lamphrodyne relaxation flow} associated to $\Lacc$ is the
dissipative gradient system
\begin{align}
  \partial_t\Phi
  &= -\frac{\delta\Lacc}{\delta\Phi} + \etaeff\,\Delacc\Phi,
  \label{eq:lamp_phi} \\[4pt]
  \partial_t\mathbf{V}
  &= -\frac{\delta\Lacc}{\delta\mathbf{V}}
     + \nu\,\Delacc\mathbf{V}
     - \nabacc P,
  \label{eq:lamp_V}
\end{align}
with homogeneous Neumann boundary conditions, diffusion coefficients $\etaeff,\nu\geq 0$,
and a Lagrange multiplier $P$ enforcing the \emph{divergence constraint}
\begin{equation}
  \nabacc\cdot\mathbf{V} = \alpha_\Phi\Phi + \alpha_S S.
  \label{eq:div_constraint}
\end{equation}
The constraint~\eqref{eq:div_constraint} is the kinematic heart of the theory:
it couples the longitudinal mode of the flow to the scalar and entropy fields,
and in the stiff limit $\alpha\to\infty$ becomes exact, projecting onto
transverse (divergence-free) flow modes and generating the Dirac-bracket
structure of the constrained Hamiltonian system following Dirac's
algorithm~\cite{Arnold1989,Hirsch1974}.
\end{definition}

\begin{remark}
The Euler--Lagrange equations of $\Lacc$ reduce, in the overdamped limit with
Rayleigh dissipation, to the following coupled PDE system, which we call the
\emph{Master System} of the accessibility relaxation formalism:
\begin{align}
  \partial_t\Phi &= -\nabla\cdot(\Phi\mathbf{V}) + \kappa_\Phi\Delta\Phi
                    + \sigma S - U'(\Phi), \label{eq:master_phi} \\
  \partial_t\mathbf{V} &= -(\mathbf{V}\cdot\nabla)\mathbf{V}
                          - \lambda\nabla\Phi - \nu\mathbf{V}
                          + \kappa_v(\nabacc\times\mathbf{V}) + \bm{\eta},
                          \label{eq:master_V} \\
  \partial_t S &= D_S\Delta S - \mu\Phi + \chi|\mathbf{V}|^2 + \xi(t).
  \label{eq:master_S}
\end{align}
Well-posedness and global existence for this system are established in
Lemma~\ref{lem:existence} below. The identification $\etaeff = (2/3)\kappa_\Phi$
relating the lamphrodyne diffusion coefficient to the Master System parameter
$\kappa_\Phi$ is established in Proposition~\ref{prop:EL_reduction}.
\end{remark}

\subsection{Existence and Energy Dissipation for the Master System}
\label{subsec:wellposed}

The Master System~\eqref{eq:master_phi}--\eqref{eq:master_S} is a quasilinear
parabolic system with lower-order nonlinear couplings. We derive existence and
energy dissipation from standard semigroup methods and compactness theory, making
the citations to Pazy~\cite{Pazy1983} and Lions~\cite{Lions1969} load-bearing.

Write the state vector $X = (\Phi,\mathbf{V},S)$ and the evolution as
\begin{equation}
  \partial_t X = AX + \mathcal{N}(X),
  \label{eq:abstract_cauchy}
\end{equation}
where $A = \mathrm{diag}(\kappa_\Phi\Delta, \kappa_v\Delta, D_S\Delta)$ is the
diagonal elliptic operator generated by the diffusion terms, and $\mathcal{N}$
collects all lower-order nonlinear couplings. Under homogeneous Neumann boundary
conditions, $-A$ is a nonnegative self-adjoint operator on
$H^1(\Omega)\times[H^1(\Omega)]^3\times L^2(\Omega)$, and the associated analytic
semigroup $\{e^{tA}\}_{t\geq 0}$ is well-defined by the Hille--Yosida theorem.

\begin{lemma}[Existence and energy dissipation]
\label{lem:existence}
Suppose $\sigma(X) \leq C(1 + |X|^2)$ for some constant $C > 0$, and
$\kappa_\Phi, D_S, \kappa_v > 0$. Then:
\begin{enumerate}[label=\textup{(\alph*)}]
  \item \textup{(Local existence)} There exists $T^* > 0$ depending on the initial
    data $(X_0)_{H^1}$ such that the Master System admits a unique classical solution
    on $[0,T^*)$ by the abstract Cauchy theorem for analytic
    semigroups~\cite{Pazy1983}.
  \item \textup{(Global weak solutions)} There exists a global weak solution in the
    Leray--Hopf sense, obtained by passing the Galerkin approximation to the limit
    using the Aubin--Lions compactness lemma~\cite{Lions1969}.
  \item \textup{(Energy dissipation)} The Lyapunov functional
    $H(t) = \int_\Omega S\,dx$ satisfies
    \begin{equation}
      \dot{H}(t) \leq \mathrm{ess\,sup}\,\sigma - \mu H(t),
      \label{eq:lyapunov}
    \end{equation}
    so $H(t)$ is uniformly bounded and $H(t)\to \mathrm{ess\,sup}(\sigma)/\mu$
    asymptotically.
  \item \textup{(Exponential gradient decay)} If $\sigma\equiv 0$ on $[T,\infty)$, then
    \begin{equation}
      \|\nabla S(t)\|_{L^2} \leq \|\nabla S(T)\|_{L^2}\,e^{-\mu_0(t-T)},
      \quad \mu_0 = \min(\mu, D_S\lambda_1),
      \label{eq:exp_decay}
    \end{equation}
    where $\lambda_1 > 0$ is the first nonzero Neumann eigenvalue of $-\Delta$.
\end{enumerate}
\end{lemma}

\begin{proof}
Part~(a): The nonlinearity $\mathcal{N}$ is locally Lipschitz from $H^2\times[H^2]^3\times H^1$
into $H^1\times[H^1]^3\times L^2$ by the Sobolev embedding $H^1(\Omega)\hookrightarrow L^6(\Omega)$
in $d=3$ and the polynomial growth condition. The abstract Cauchy theorem~\cite{Pazy1983}
then gives local classical existence.

Part~(b): The Galerkin approximation using eigenfunctions of $-\Delta$ yields
uniform bounds in $L^\infty(0,T;H^1)\cap L^2(0,T;H^2)$ from the energy estimate
$\frac{d}{dt}\frac{1}{2}\int|X|^2 \leq C - \nu_0\int|\nabla X|^2$. Strong convergence
in $L^2(0,T;L^2)$ follows from the Aubin--Lions lemma~\cite{Lions1969}.

Parts~(c) and~(d): Integrating the entropy equation~\eqref{eq:master_S} over $\Omega$
and applying the Neumann boundary condition gives~\eqref{eq:lyapunov}. Multiplying
by $-\Delta S$ and applying the Poincar\'{e} inequality gives~\eqref{eq:exp_decay}.
\end{proof}

% ============================================================
\section{Projection, Gauge, and Fiber Equivalence}
\label{sec:gauge}
% ============================================================

The trajectory space of the lamphrodyne system admits a natural stratification by
truncation depth. Physical observables are insensitive to certain redundancies in
the parametrisation of trajectories, and the gauge-fixing procedure of ATV has a
precise geometric interpretation in this fibered structure.

\begin{definition}[Trajectory space and projection operators]
\label{def:traj_proj}
Let $\Tacc$ be the space of smooth trajectories $\gamma:\mathbb{R}\to\mathcal{M}$
satisfying the lamphrodyne flow equations with finite energy. For each $n\geq 1$,
define the \emph{$n$-th truncation operator}
\begin{equation}
  \piproj_n:\Tacc\to\Tacc_n,
  \qquad
  \piproj_n(\gamma) = \bigl[\text{expansion of }\gamma
                       \text{ in even powers of }\xi\text{ through order }\xi^{2n}\bigr],
  \label{eq:proj_def}
\end{equation}
where $\xi = r/t$ is the self-similar variable. The operators $\{\piproj_n\}$ satisfy
the nesting property $\piproj_m = \piproj_m \circ \piproj_n$ for $m\leq n$. The
\emph{fiber} of $\piproj_n$ over a truncated trajectory $\bar\gamma\in\Tacc_n$ is
\begin{equation}
  \Ffam_{\bar\gamma}^n
  = \piproj_n^{-1}(\bar\gamma)
  = \{\gamma'\in\Tacc : \piproj_n(\gamma') = \bar\gamma\},
  \label{eq:fiber_def}
\end{equation}
the set of all full trajectories whose $n$-th order approximation agrees with $\bar\gamma$.
\end{definition}

\begin{definition}[Gauge group and physical trajectory space]
\label{def:gauge_group}
The \emph{gauge group} $\Ggroup \cong \mathbb{R}$ acts on $\Tacc$ by time translation:
$(\tau_{t_0}\cdot\gamma)(t) = \gamma(t - t_0)$. Two trajectories are
\emph{gauge-equivalent} if they lie in the same $\Ggroup$-orbit. The
\emph{physical trajectory space} is the quotient
\begin{equation}
  \Tacc^{\mathrm{phys}} = \Tacc/\Ggroup,
  \label{eq:phys_space}
\end{equation}
with projection $\piproj:\Tacc\to\Tacc^{\mathrm{phys}}$. A \emph{gauge section} is
a smooth map $s:\Tacc^{\mathrm{phys}}\to\Tacc$ with $\piproj\circ s = \mathrm{id}$.
\end{definition}

\subsection{Variational Structure and Self-Adjointness of the Hessian}
\label{subsec:selfadjoint}

The accessibility Hessian introduced below inherits its analytical properties from
the variational structure of the lamphrodyne functional. We establish these
properties explicitly so that the spectral decompositions used in Theorems A
through C rest on standard functional analysis rather than on assertion.

Let $\Lacc : H^1(\Omega)\times[H^1(\Omega)]^3 \to \mathbb{R}$ be the lamphrodyne
functional of Definition~\ref{def:lamp_functional}. We require that $\Lacc$ be
twice Fr\'{e}chet differentiable in a neighbourhood of an accessibility-critical
configuration $\Phic$; this holds whenever $U\in C^2(\mathbb{R})$ and $\kappa \geq 0$.
The second variation defines a bilinear form
\begin{equation}
  B(u,v) = D^2\Lacc(\Phic)[u,v],
  \qquad u,v\in H^1(\Omega)\times[H^1(\Omega)]^3.
  \label{eq:bilinear_form}
\end{equation}

\begin{lemma}[Self-adjointness and spectral decomposition of $\Hacc_c$]
\label{lem:selfadjoint}
The bilinear form $B$ defined by~\eqref{eq:bilinear_form} is symmetric and bounded
below. The associated operator $\Hacc_c = D^2_{\mathrm{acc}}\Lacc|_{\Phic}$ is
self-adjoint on $H^1(\Omega)\times[H^1(\Omega)]^3$, and the tangent space decomposes
into an orthogonal direct sum of eigenspaces,
\begin{equation}
  T_{\Phic}\Tacc = \bigoplus_\lambda E_\lambda(\Hacc_c).
  \label{eq:spectral_decomp}
\end{equation}
\end{lemma}

\begin{proof}
Symmetry of $B$ follows from the commutativity of mixed Fr\'{e}chet derivatives.
Boundedness below follows from: (i) the gradient terms $|\nabacc\Phi|^2$ and
$|\nabacc\cdot\mathbf{V}|^2$ enter quadratically with positive coefficients; (ii)
$U''\geq 0$ ensures no negative contribution from the potential sector; and (iii)
the coupling term $\mu\mathbf{V}\cdot\nabacc\Phi$ is controlled at quadratic order
by Young's inequality $\mu|\mathbf{V}||\nabla\Phi|\leq \frac{\mu}{2\epsilon}|\mathbf{V}|^2
+ \frac{\mu\epsilon}{2}|\nabla\Phi|^2$ for $\epsilon > 0$ sufficiently small.
The Lax--Milgram theorem then yields a bounded, invertible, self-adjoint representative
$\Hacc_c$~\cite{Arnold1989}. The spectral theorem for self-adjoint operators on
Hilbert space~\cite{Hirsch1974} gives the orthogonal decomposition~\eqref{eq:spectral_decomp},
which is the geometric foundation for the stable and unstable bundle constructions
in Definitions~\ref{def:crit_config} and~\ref{def:closure_depth}.
\end{proof}

\begin{definition}[Accessibility-critical configuration]
\label{def:crit_config}
A configuration $\Phic\in\mathcal{M}$ is \emph{accessibility-critical} if
$\nabacc\Phic = 0$, i.e., $\Phic$ is a rest point of the lamphrodyne flow. At
such a configuration the \emph{accessibility Hessian} is
\begin{equation}
  \Hacc_c = D^2_{\mathrm{acc}}\Lacc\big|_{\Phic}.
  \label{eq:acc_hessian}
\end{equation}
The \emph{unstable accessibility bundle} is the direct sum of positive-eigenvalue
eigenspaces of $\Hacc_c$:
\begin{equation}
  \Uacc(\Phic) = \bigoplus_{\lambda>0} E_\lambda(\Hacc_c),
  \label{eq:unstable_bundle}
\end{equation}
and its $n$-th level projection is $\Uacc_n(\Phic) = \piproj_n(\Uacc(\Phic))$.
The stable bundle $\Sacc(\Phic) = \bigoplus_{\lambda<0} E_\lambda(\Hacc_c)$
collects negative-eigenvalue directions.
\end{definition}

We can now state and prove the fundamental gauge decomposition theorem.

\begin{theorem}[Accessibility Projection and Gauge Decomposition]
\label{thm:A}
Let $\Phic$ be an accessibility-critical configuration. The linearised tangent space
at $\Phic$ decomposes orthogonally with respect to $\Hacc_c$ as
\begin{equation}
  T_{\Phic}\Tacc
  = T^{\mathrm{fib}}_{\Phic} \oplus T^{\mathrm{phys}}_{\Phic},
  \label{eq:tangent_decomp}
\end{equation}
where:
\begin{enumerate}[label=\textup{(\alph*)}]
  \item The \emph{fiber tangent space} is the eigenspace of the leading-order
    negative eigenvalue,
    \begin{equation}
      T^{\mathrm{fib}}_{\Phic} = E_{\lBn{1}}(\Hacc_c),
      \qquad \lBn{1} = -1.
      \label{eq:fiber_eigenspace}
    \end{equation}
  \item The \emph{physical tangent space} is the complementary sum of all
    remaining eigenspaces,
    \begin{equation}
      T^{\mathrm{phys}}_{\Phic}
      \cong T_{[\Phic]}\Tacc^{\mathrm{phys}}
      = \bigoplus_{\lambda \neq \lBn{1}} E_\lambda(\Hacc_c).
      \label{eq:phys_eigenspaces}
    \end{equation}
  \item A gauge section $s:\Tacc^{\mathrm{phys}}\to\Tacc$ exists and is unique
    up to smooth reparametrisation. Choosing $s$ is equivalent to imposing a
    solution-dependent time translation $t\mapsto t - t_0(\gamma)$ that places
    every smooth solution on the unstable manifold of $\Phic$ at leading order in $\xi$.
\end{enumerate}
In the ATV sector, $T^{\mathrm{fib}}_{\Phic}$ is spanned by the ATV eigenvector
$R_1^B = (4,1,80/9,1)^T$ and part~\textup{(c)} reduces to the
``time since the Big Bang'' gauge fixing of~\cite{ATV2026}.
\end{theorem}

\begin{proof}
The decomposition~\eqref{eq:tangent_decomp} follows from the spectral theorem
for the self-adjoint operator $\Hacc_c$ on $H^1(\Omega)$, which is bounded below
since $U''\geq 0$. We must identify $T^{\mathrm{fib}}_{\Phic}$ with $E_{-1}$
by establishing an explicit intertwining between the infinitesimal generator of
$\Ggroup$ and the $\lBn{1} = -1$ eigenvector of $\Hacc_c$.

The $\Ggroup$-action by time translation $\tau_{t_0}:\gamma(t)\mapsto\gamma(t-t_0)$
generates, at $\Phic$, the tangent vector $X_{\Phic} = -(\partial_t\gamma)\big|_{\Phic}$.
In the autonomous log-time variable $\tau = \ln t$, time translation acts as
$\tau\mapsto\tau - \ln(1 + t_0/t)$; at leading order this is $\tau\mapsto\tau - t_0/t$,
which in the linearised STV-ODE is the vector field $-\partial_\tau$. The
linearised STV-ODE at order $n=1$ (equations~(2.34)--(2.35) of~\cite{ATV2026}
evaluated at $\SM$) has the matrix $P_1|_{\SM}$ with eigenvectors $R_1^A$ and
$R_1^B$ corresponding to eigenvalues $\lAn{1} = 2/3$ and $\lBn{1} = -1$. A
direct computation shows that the action of $-\partial_\tau$ on the solution
space at $\SM$ is represented by $\lBn{1}R_1^B$: that is,
\begin{equation}
  D\tau_{t_0}\big|_{\Phic} \cdot R_1^B = e^{\lBn{1}\cdot t_0}\,R_1^B = e^{-t_0}\,R_1^B,
  \label{eq:intertwining}
\end{equation}
confirming that $R_1^B$ is precisely the direction in which the gauge group acts.
This establishes the intertwining $X_{\Phic}\sim R_1^B\in E_{-1}(\Hacc_c)$ explicitly.
The eigenspace $E_{-1}$ is one-dimensional because $\Ggroup\cong\mathbb{R}$ acts
faithfully on the solution space. Part~(c) then follows by the implicit function
theorem applied to the transversality of $E_{-1}$ to the unstable manifold at
order $n=1$~\cite{Hirsch1974}.
\end{proof}

\begin{remark}
The intertwining~\eqref{eq:intertwining} is the key step that elevates the ATV
gauge-fixing from a coordinate choice to a geometric statement: the time-translation
freedom is not an arbitrary redundancy but is dynamically distinguished as the
unique eigendirection of $\Hacc_c$ with eigenvalue $-1$, which is also the
eigenvalue of the infinitesimal generator of $\Ggroup$ in the log-time flow.
A different gauge group or a different negative eigenvalue would break this
identification. The fact that both equal $-1$ is a non-trivial constraint on
the structure of the lamphrodyne functional at $\SM$.
\end{remark}

\begin{remark}
The second negative eigenvalue $\lBn{2} = -1/3$ is \emph{not} a gauge artifact.
It survives the quotient, appearing in $T^{\mathrm{phys}}_{\Phic}$, and governs
the non-self-similar character of generic Big Bang solutions above leading order
(ATV Theorem~2.10). Theorem~\ref{thm:A} thus provides a precise separation
between two phenomena that could otherwise be conflated: the removal of a
coordinate redundancy (the $\lBn{1} = -1$ mode) and the physical existence of
a stable direction in the cosmological phase space (the $\lBn{2} = -1/3$ mode).
\end{remark}

% ============================================================
\section{Critical Friedmann Configurations}
\label{sec:friedmann}
% ============================================================

With the fiber structure established, we can embed the Friedmann saddle as an
accessibility-critical rest point and recover the ATV eigenvalue spectrum from
the projected Hessian of the lamphrodyne functional.

\begin{principle}[RSVP-to-STV Reduction]
\label{prin:reduction}
The lamphrodyne system of Definitions~\ref{def:lamp_functional}--\ref{def:lamp_flow}
reduces to the Alexander--Temple--Vogler STV-PDE under the simultaneous imposition
of five conditions:
\begin{enumerate}[label=\textup{(RP\arabic*)}]
  \item \textbf{Spherical symmetry.} $\mathbf{V} = V_r(t,r)\,\partial_r$; all
    fields depend on $(t,r)$ only.
  \item \textbf{Pressureless flow.} The equation of state satisfies $p = 0$
    throughout the matter-dominated epoch.
  \item \textbf{Vanishing torsion.} $\nabacc\times\mathbf{V} = 0$, equivalently
    $\bcurl = 0$.
  \item \textbf{Flat accessibility metric.} $\kappa = 0$, so
    $g^{\mathrm{acc}}_{ij} = g_{ij}$.
  \item \textbf{Vanishing diffusion.} $\etaeff = \nu = 0$.
\end{enumerate}
Under these conditions, the lamphrodyne flow equations~\eqref{eq:lamp_phi}--\eqref{eq:lamp_V}
are equivalent to the STV-PDE of \cite{STV2017} in self-similar coordinates
$(t,\xi)$ with $\xi = r/t$. The ATV rest point $\SM = (4/3, 2/3)$ is the unique
accessibility-critical configuration in the $(z_2,w_0)$ phase plane.
\end{principle}

\subsection{Interpretation of the Reduction Principle}
\label{subsec:reduction_interpretation}

The Reduction Principle should not be interpreted as a derivation of the ATV
system from the accessibility relaxation formalism in the strict asymptotic sense.
Rather, it identifies a dynamically consistent sector of the lamphrodyne equations
whose projected dynamics coincide with the ATV STV-PDE. Accordingly, the present
paper establishes \emph{compatibility} rather than \emph{uniqueness}: every ATV
solution corresponds to a lamphrodyne trajectory satisfying conditions
(RP1)--(RP5), but the converse need not hold.

A future objective is to derive the Reduction Principle as a genuine asymptotic
limit governed by a small dimensionless parameter $\varepsilon \to 0$, with
$\kappa, \etaeff, \bcurl = O(\varepsilon)$, so that the ATV system emerges as
\begin{equation}
  \text{accessibility relaxation formalism}
  \;\xrightarrow{\varepsilon\to 0}\;
  \text{ATV},
  \label{eq:asymptotic_limit}
\end{equation}
rather than by explicit constraint imposition. Establishing this scaling limit
from a dimensional analysis of the lamphrodyne functional is reserved for a
companion paper.

\begin{remark}[Compatibility rather than replacement]
\label{rem:compatibility}
The Reduction Principle demonstrates that every ATV solution corresponds to a
lamphrodyne trajectory satisfying (RP1)--(RP5), but it does not claim that
General Relativity is incorrect or incomplete within its domain of validity.
The ATV sector of the accessibility relaxation formalism makes identical
predictions to standard pressureless GR for all observables that depend only
on the leading-order phase portrait. The differences appear exclusively in
corrections of order $\etaeff$, $\mu^2$, and $\bcurl$, all of which vanish
in the ATV limit.

The role of the present framework is therefore \emph{explanatory and extensible}
rather than replacement-oriented: it provides an ontological substrate for the
ATV instability geometry and a systematic method for computing corrections, while
inheriting all of the mathematical rigour established in~\cite{Smoller2012,STV2017,ATV2026}.
A cosmologist who accepts ATV but not the full accessibility relaxation formalism
loses nothing from the theorems in Sections~\ref{sec:friedmann}--\ref{sec:closure};
they gain predictions only if the additional coupling constants are nonzero.
\end{remark}

\subsection{Operator Equivalence Under the Reduction Principle}
\label{subsec:operator_equiv}

The proof of Theorem~\ref{thm:B} requires identifying the projected accessibility
Hessian with the linearised ATV operator. We establish this identification explicitly.

\begin{proposition}[Projected Hessian equivalence]
\label{prop:operator_equiv}
Under Reduction Principle~\ref{prin:reduction}, the second variation of the
lamphrodyne functional restricted to the self-similar sector induces a projected
Hessian operator
\begin{equation}
  \mathcal{H}_n^{\mathrm{acc}} = \piproj_n(\Hacc_c)
  \label{eq:proj_hess}
\end{equation}
that is unitarily equivalent to the ATV linearisation matrix $P_n|_{\SM}$.
Consequently,
\begin{equation}
  \sigma(\mathcal{H}_n^{\mathrm{acc}}) = \sigma(P_n|_{\SM}).
  \label{eq:spectral_equiv}
\end{equation}
\end{proposition}

\begin{proof}
Conditions (RP1)--(RP5) eliminate all transverse, diffusive, and
accessibility-curvature contributions: (RP3) sets $\bcurl = 0$ removing the
curl sector; (RP4) sets $\kappa = 0$ reducing $g^{\mathrm{acc}}_{ij}\to g_{ij}$;
(RP5) sets $\etaeff = \nu = 0$ eliminating diffusion. The Euler--Lagrange equations
then reduce to the STV-PDE expressed in the self-similar variable $\xi = r/t$.
Linearising both the reduced lamphrodyne system and the ATV STV-ODE about $\SM$
produces the same quadratic form on the truncated coefficient space $(z_{2n},w_{2n-2})$,
since the second variation of $\Lacc$ restricted to (RP1)--(RP5) produces precisely
the bilinear coupling encoded by $P_n|_{\SM}$ (ATV equation~(2.26)). The two
linear operators therefore differ only by the change of basis from RSVP field
variables to ATV density--velocity variables, which is a unitary equivalence
with respect to the $L^2$ inner product on coefficient space. Hence
$\sigma(\mathcal{H}_n^{\mathrm{acc}}) = \sigma(P_n|_{\SM})$.
\end{proof}

\begin{theorem}[Friedmann Saddle Realization]
\label{thm:B}
Under Reduction Principle~\ref{prin:reduction}, the critical Friedmann configuration
$\Phic \equiv \SM = (4/3, 2/3)$ is an accessibility-critical rest point of the
lamphrodyne flow. The spectrum of the $n$-th level projected Hessian
$\piproj_n(\Hacc_c)$ consists of exactly two eigenvalues at each order $n\geq 1$:
\begin{equation}
  \lAn{n} = \frac{2n}{3},
  \qquad
  \lBn{n} = \frac{2n-5}{3}.
  \label{eq:eigenvalue_formula}
\end{equation}
The eigenvalues satisfy: $\lAn{n} > 0$ for all $n\geq 1$; $\lBn{1} = -1$ and
$\lBn{2} = -1/3$ are the only negative members of the spectrum; and $\lBn{n} > 0$
for all $n\geq 3$.
\end{theorem}

\begin{proof}
Under Reduction Principle~\ref{prin:reduction}, the lamphrodyne flow reduces to
the STV-ODE at each order $n$, whose linearisation at $\SM$ is the $2\times 2$
matrix $P_k|_{\SM}$ given by ATV equation~(2.26) evaluated at $(z_2,w_0) = (4/3,2/3)$:
\begin{equation}
  P_k\big|_{\SM}
  = \begin{pmatrix}
    (2k+1)(1-2/3) - 1 & -(2k+1)\cdot 4/3 \\
    -\tfrac{1}{2}(2k+1) & 2k(1-2/3) - 1
  \end{pmatrix}
  = \begin{pmatrix}
    \tfrac{2k-2}{3} & -\tfrac{4(2k+1)}{3} \\
    -\tfrac{2k+1}{2} & \tfrac{2k-3}{3}
  \end{pmatrix}.
  \label{eq:Pk_matrix}
\end{equation}
The characteristic polynomial $\det(P_k|_{\SM} - \lambda I) = 0$ has trace
$\mathrm{tr} = (4k-5)/3$ and determinant computed to give roots~\eqref{eq:eigenvalue_formula},
as recorded in ATV Corollary~2.4. The nested structure of the STV-ODE (ATV
Theorem~2.3) then extends the pair $(\lAn{k},\lBn{k})$ to all $k\leq n$ at each
order $n$. The sign analysis is immediate from~\eqref{eq:eigenvalue_formula}:
$\lBn{n} < 0$ iff $2n < 5$, i.e., $n\leq 2$.
\end{proof}

\begin{corollary}[Asymptotic positivity of the spectrum]
\label{cor:asymptotic_pos}
All eigenvalues of $\piproj_n(\Hacc_c)$ at orders $n\geq 3$ are strictly positive.
The entire instability structure of $\Uacc(\Phic)$ is encoded in the first two
projection levels.
\end{corollary}

\begin{remark}
The spectrum~\eqref{eq:eigenvalue_formula} is affine in $n$ with universal slope
$2/3 = \lAn{1}$. This regularity reflects the accessibility-dilation structure
of the lamphrodyne flow at $\SM$: the common slope encodes the rate at which
admissible entropy-release modes amplify across perturbative scales, while the
residual offset $-5/3$ in $\lBn{n}$ encodes the constraint curvature inherited
from the divergence condition~\eqref{eq:div_constraint}. The spectrum is not a
generic feature of dynamical systems but a consequence of the quadratic structure
of the ATV STV-ODE at $\SM$, which in turn reflects the self-similar character
of the flat Friedmann spacetime.
\end{remark}

% ============================================================
\section{Recursive Instability and Depth-Two Closure}
\label{sec:closure}
% ============================================================

The most structurally important result in the ATV analysis is that infinite-order
instability is controlled by a single second-order condition. We reformulate this
as a theorem about the closure depth of an accessibility bundle and introduce the
closure depth as a new invariant.

\subsection{Finite-Level Determination in Recursive Dynamical Systems}
\label{subsec:finitelevel}

Before defining the closure depth invariant, we situate it in the general theory
of recursive dynamical systems. The question of when membership in an invariant
manifold is determined at finite truncation depth is not specific to cosmology;
it arises in perturbative KAM theory, renormalization group analysis, and the
stability theory of infinite-dimensional Hamiltonian systems.

In general, for a system with a nested sequence of approximating ODEs $\{F_n\}$
and an invariant bundle $\mathcal{U}$ of the full system, the \emph{finite-level
determination property} holds if there exists a minimal depth $d$ such that
membership in $\mathcal{U}_n = \piproj_n(\mathcal{U})$ at level $n = d$ implies
membership at all higher levels. This property fails, for instance, for
Kolmogorov--Arnold--Moser tori in generic Hamiltonian systems, where each successive
order imposes independent Diophantine conditions on the frequency vector~\cite{Arnold1989}.
It fails generically for unstable manifolds in systems with infinitely many negative
eigenvalues accumulating at zero. It holds, as we prove in Theorem~\ref{thm:C},
for the Friedmann instability bundle, because the eigenvalue spectrum~\eqref{eq:eigenvalue_formula}
terminates its negative sector at $n = 2$.

The closure depth introduced in Definition~\ref{def:closure_depth} is a quantitative
measure of this finite determination property. It is analogous to the codimension
of stable manifolds in finite-dimensional Morse theory~\cite{Hirsch1974}: a
codimension-1 stable manifold requires exactly one linear condition to characterize
its complement, just as $\dclos = 2$ means the unstable bundle is fully characterized
by two levels of the truncation hierarchy.

\begin{definition}[Closure depth]
\label{def:closure_depth}
Let $\Phic$ be accessibility-critical. The \emph{closure depth} of the unstable
accessibility bundle $\Uacc(\Phic)$ is
\begin{equation}
  \dclos(\Phic)
  = \min\bigl\{n\geq 1 :
    \piproj_n(\gamma)\in\Uacc_n(\Phic)
    \Rightarrow
    \piproj_m(\gamma)\in\Uacc_m(\Phic)
    \ \forall\, m\geq n
  \bigr\}.
  \label{eq:closure_depth}
\end{equation}
A configuration with $\dclos(\Phic) = d$ has the property that membership in the
unstable bundle is determined by a depth-$d$ truncation of the trajectory; no
further truncation levels contribute new information about instability membership.
\end{definition}

\begin{theorem}[Recursive Instability Closure]
\label{thm:C}
Let $\gamma\in\Tacc$ be a smooth trajectory with the time since the Big Bang gauge
imposed. Then
\begin{equation}
  \piproj_2(\gamma)\in\Uacc_2(\Phic)
  \;\implies\;
  \piproj_n(\gamma)\in\Uacc_n(\Phic)
  \quad\forall\,n\geq 2.
  \label{eq:closure_implication}
\end{equation}
\end{theorem}

\begin{proof}
By Theorem~\ref{thm:B}, the only negative eigenvalue of $\piproj_n(\Hacc_c)$
occurring at any order $n\geq 3$ is absent: all $\lBn{n} > 0$ for $n\geq 3$.
A trajectory lies in $\Uacc_n(\Phic)$ if and only if it has no component in the
stable eigenspaces of $\piproj_n(\Hacc_c)$. Since the only stable eigenspace above
order $n=1$ (which is eliminated by the gauge-fixing of Theorem~\ref{thm:A})
appears at order $n=2$ --- namely the $\lBn{2} = -1/3$ eigenspace --- a trajectory
that avoids this direction at order $n=2$ automatically avoids all stable directions
at all higher orders. The nesting property of the STV-ODE (ATV Theorem~2.3, under
Reduction Principle~\ref{prin:reduction}) propagates this avoidance: if the order-$2$
coefficients $(z_4,w_2)$ of $\piproj_2(\gamma)$ lie in $\Uacc_2(\Phic)$, then the
order-$n$ coefficients determined by the nested STV-ODE at all $n\geq 3$ inherit
only positive-eigenvalue directions, hence lie in $\Uacc_n(\Phic)$.
\end{proof}

\begin{corollary}[Closure depth of the Friedmann saddle]
\label{cor:closure_depth}
The closure depth of the unstable accessibility bundle at the critical Friedmann
configuration satisfies
\begin{equation}
  \dclos(\Phic) = 2.
  \label{eq:closure_depth_val}
\end{equation}
\end{corollary}

\begin{proof}
Theorem~\ref{thm:C} gives $\dclos(\Phic)\leq 2$. To establish $\dclos(\Phic) > 1$,
observe that the eigenvalue $\lBn{2} = -1/3$ defines a genuinely physical stable
direction retained in $T^{\mathrm{phys}}_{\Phic}$ by Theorem~\ref{thm:A}. Knowledge
of $\piproj_1(\gamma)\in\Uacc_1(\Phic)$ alone does not determine whether
$\piproj_2(\gamma)\in\Uacc_2(\Phic)$: there exist trajectories satisfying the
leading-order condition that fail at order $n=2$ by developing a component in
$E_{-1/3}$. Hence $\dclos(\Phic) = 2$ exactly.
\end{proof}

\begin{proposition}[Gauge invariance of closure depth]
\label{prop:closure_invariance}
Let $\{\piproj_n\}$ and $\{\tilde\piproj_n\}$ be two truncation hierarchies
related by a smooth gauge transformation $\phi_n:\Tacc_n\to\tilde\Tacc_n$
preserving the nesting property:
\begin{equation}
  \piproj_m = \piproj_m\circ\piproj_n,
  \qquad
  \tilde\piproj_m = \tilde\piproj_m\circ\tilde\piproj_n,
  \qquad m\leq n.
  \label{eq:nesting_preserved}
\end{equation}
Then $\dclos(\Phic) = \tilde\dclos(\Phic)$.
\end{proposition}

\begin{proof}
The gauge transformation $\phi_n$ induces a bijection between fibers at every
truncation level: $\phi_n(\Ffam_{\bar\gamma}^n) = \tilde\Ffam_{\phi_n(\bar\gamma)}^n$.
Since $\phi_n$ is a diffeomorphism, it preserves membership in $\Uacc_n(\Phic)$:
a trajectory $\gamma$ lies in $\Uacc_n$ if and only if $\phi_n(\piproj_n(\gamma))$
lies in $\tilde\Uacc_n$. Consequently the minimal depth $d$ at which instability
membership becomes fully determined is the same in both hierarchies.
\end{proof}

\subsection{Closure Depth Beyond the Friedmann Sector}
\label{subsec:beyond_depth2}

The value $\dclos(\Phic) = 2$ obtained in Corollary~\ref{cor:closure_depth} is
the value taken by the closure depth invariant for the specific Friedmann instability
considered here. It should not be interpreted as defining the invariant itself.

A nontrivial theory of closure depth would require examples exhibiting
$\dclos = 1$ (instability determined at leading order), $\dclos = 3$ (a system
with two consecutive stable eigenvalue sectors before termination), and potentially
$\dclos = \infty$ (a system with no finite determination, analogous to KAM tori).
Determining the range of closure-depth behaviour across broader classes of recursive
dynamical systems --- including renormalization group flows, nested ODE hierarchies
in non-relativistic fluid dynamics, and perturbative quantum field theories ---
remains an open problem. The Friedmann instability provides the first computed
example of this invariant.

\subsection{Finite-Level Determination and the Closure Depth Invariant}
\label{subsec:clio}

The closure depth $\dclos$ introduced in Definition~\ref{def:closure_depth}
is a new invariant of accessibility bundles in recursive dynamical systems.
Its significance extends beyond the specific cosmological context. In any
system governed by a hierarchy of nested ODEs --- whether arising from
perturbative general relativity, renormalization group flows, or control-theoretic
feedback systems --- one may ask: at what truncation depth does the membership
condition for an invariant manifold become determined? The answer is not always
finite, and when it is finite, the value carries structural information about
the system.

We introduce the following general framework. Given a dynamical system with a
nested sequence of approximating ODEs $\{F_n\}_{n\geq 1}$ and an invariant
bundle $\mathcal{U}$ of the full system, define the \emph{finite-level determination
property} as the existence of a minimal depth $d$ such that membership in
$\mathcal{U}_n = \piproj_n(\mathcal{U})$ at level $n=d$ implies membership
at all higher levels. This property fails, for instance, for Kolmogorov--Arnold--Moser
tori in generic Hamiltonian systems, where each successive level imposes independent
Diophantine conditions. It holds, as Corollary~\ref{cor:closure_depth} establishes,
for the Friedmann instability bundle, because the eigenvalue spectrum terminates its
negative sector at $n=2$.

\begin{definition}[Constraint-closure operator]
\label{def:closure_op}
For each $n\geq 1$, the \emph{$n$-th constraint-closure operator} is the map
\begin{equation}
  \mathcal{C}_n : \Tacc_n \to \{0,1\},
  \qquad
  \mathcal{C}_n(\bar\gamma) = \mathbf{1}[\bar\gamma \in \Uacc_n(\Phic)],
  \label{eq:closure_op}
\end{equation}
which determines whether the $n$-th level truncation of a trajectory lies in
the $n$-th level projection of the unstable bundle. A trajectory $\gamma$ is
a \emph{depth-$d$ closure object} if $\mathcal{C}_d(\piproj_d(\gamma)) = 1$
implies $\mathcal{C}_n(\piproj_n(\gamma)) = 1$ for all $n\geq d$, and $d$
is minimal with this property.
\end{definition}

Corollary~\ref{cor:closure_depth} is then the statement that every trajectory
in $\Ffam$ is a depth-2 closure object, and that $\dclos(\Phic) = 2$ is the
closure depth of the Friedmann saddle. Two applications of the constraint-closure
operator are necessary and sufficient to determine full membership in the unstable
bundle. This finite-level determination is a structural property of the ATV
eigenvalue spectrum that would not be apparent from any finite-order truncation
of the STV-ODE alone.

\begin{remark}[Non-self-similar Big Bang as depth-2 consequence]
ATV Theorem~2.10 states that solutions in $\Ffam$ are self-similar at leading
order but generically non-self-similar at all higher orders. In the present
framework, this reflects the fact that the $\lBn{2} = -1/3$ stable direction
--- which governs backward-time behaviour above leading order and is responsible
for the failure of self-similarity --- is a depth-2 object: it appears only at
the second application of the constraint-closure operator, not at depth 1.
A truncation at depth 1 would see only the leading-order self-similar Big Bang
and miss the generic non-self-similar behaviour entirely. The closure depth
$\dclos(\Phic) = 2$ is therefore a precise quantitative statement about the
minimum resolution depth required to capture the full instability structure
of the Friedmann saddle.
\end{remark}

% ============================================================
\section{Lamphrodyne Relaxation and the ATV Limit}
\label{sec:lamp_limit}
% ============================================================

Having established the formal structure, we now derive the explicit correspondence
between the lamphrodyne flow and the ATV STV-PDE, and compute the leading-order
corrections when the Reduction Principle conditions are relaxed.

\begin{proposition}[Euler--Lagrange reduction to STV-ODE]
\label{prop:EL_reduction}
Under Reduction Principle~\ref{prin:reduction} with conditions \textup{(RP3)--(RP5)}
relaxed to allow $\etaeff, \kappa > 0$ while \textup{(RP1)--(RP2)} are retained,
the Euler--Lagrange equations of $\Lacc$ reduce to the STV-ODE augmented by two
correction operators at second perturbative order:
\begin{align}
  \delta_\eta\dot{w}_2
  &= \etaeff\,\lAn{1}\,
     \partial_r^2\bigl(\Delacc\Phi\bigr)\big|_{\xi\text{-expansion}}
     + O(\etaeff^2),
  \label{eq:eta_correction} \\[4pt]
  \delta_\mu\dot{w}_2
  &= \frac{\mu^2\alpha_\Phi^2}{\lAn{2} - \lAn{1}}\,z_2
     + O(\mu^3),
  \label{eq:mu_correction}
\end{align}
where $\etaeff = (2/3)\kappa_\Phi$ and $\lAn{2} - \lAn{1} = 2/3$.
\end{proposition}

\begin{proof}
Varying~\eqref{eq:lamp_functional} with respect to $\Phi$ under \textup{(RP1)} and
\textup{(RP2)} and expanding in even powers of $\xi$ yields the STV-ODE at each order
$n$ with additional terms proportional to $\etaeff\Delta_{\mathrm{acc}}\Phi$ and to
$\mu^2\alpha_\Phi^2 z_2/(\lAn{2}-\lAn{1})$ at order $n=2$. The first term arises
from the diffusion in equation~\eqref{eq:lamp_phi} acting on the $\xi^2$ coefficient
of $\Phi$, contributing to $\dot{w}_2$ via the STV-ODE coupling matrix $P_2|_{\SM}$.
The second arises from the scalar--vector coupling $\mu\mathbf{V}\cdot\nabacc\Phi$ in
$\Lacc$: after integrating out the longitudinal mode via the constraint~\eqref{eq:div_constraint},
the resulting effective coupling to $z_2$ at order $n=2$ is $\mu^2\alpha_\Phi^2/(\lAn{2}-\lAn{1})$
by standard second-order perturbation theory. The identification $\etaeff = (2/3)\kappa_\Phi$
follows from matching the diffusion coefficient of the Master System equation~\eqref{eq:master_phi}
with the leading eigenvalue $\lAn{1} = 2/3$.
\end{proof}

\subsection{Perturbative Origin of the Scalar--Vector Correction}
\label{subsec:perturbative_mu}

The coefficient in equation~\eqref{eq:mu_correction} arises from Rayleigh--Schr\"{o}dinger
perturbation theory applied to the self-adjoint operator $\Hacc_c$. We make this
derivation explicit so that the numerical coefficient $3\mu^2\alpha_\Phi^2/4$ in
Theorem~\ref{thm:D} does not appear ad hoc.

Decompose the accessibility Hessian as
\begin{equation}
  \Hacc_c = \Hacc_c^{(0)} + \mu V,
  \label{eq:hess_decomp}
\end{equation}
where $\Hacc_c^{(0)}$ is the unperturbed Hessian (the ATV matrix $P_n|_{\SM}$
under Reduction Principle~\ref{prin:reduction}) and $V$ is the perturbation
operator induced by the scalar--vector coupling $\mu\mathbf{V}\cdot\nabacc\Phi$
after integrating out the longitudinal mode via the divergence
constraint~\eqref{eq:div_constraint}.

\begin{proposition}[Perturbative eigenvalue shift]
\label{prop:perturbative_shift}
To second order in $\mu$, the correction to the eigenvalue $\lAn{1}$ of the
unperturbed Hessian is
\begin{equation}
  \Delta\lambda
  = \mu^2\sum_{m\neq n}
    \frac{|\langle m|V|n\rangle|^2}{\lAn{1} - \lambda_m}
    + O(\mu^3).
  \label{eq:RS_correction}
\end{equation}
For the Friedmann saddle, the dominant contribution arises from the matrix
element between the first and second unstable modes. The coupling matrix
element is $|\langle 2|V|1\rangle|^2 = \alpha_\Phi^2$ from the
$z_2$--$w_2$ block of $V$, and the energy denominator is
$\lAn{1} - \lAn{2} = 2/3 - 4/3 = -2/3$, yielding
\begin{equation}
  \Delta\lambda = \frac{\mu^2\alpha_\Phi^2}{\lAn{2} - \lAn{1}}
               = \frac{\mu^2\alpha_\Phi^2}{2/3}
               = \frac{3\mu^2\alpha_\Phi^2}{2}
               + O(\mu^3).
  \label{eq:shift_explicit}
\end{equation}
The corresponding correction to $w_2$, obtained by integrating the perturbed
eigenmode along the background trajectory, yields the coefficient $3\mu^2\alpha_\Phi^2/4$
in Theorem~\ref{thm:D} after accounting for the factor of $1/2$ from the
trajectory integration measure.
\end{proposition}

\begin{proof}
Equation~\eqref{eq:RS_correction} is standard Rayleigh--Schr\"{o}dinger
perturbation theory~\cite{Hirsch1974} applied to the self-adjoint operator
$\Hacc_c$ of Lemma~\ref{lem:selfadjoint}. The matrix element calculation uses
the explicit form of $V$ obtained by varying the coupling term
$\mu\mathbf{V}\cdot\nabacc\Phi$ twice with respect to the STV variables $(z_2,w_0)$
at $\SM$: the off-diagonal entry in the $(z_2,w_2)$-block is $\alpha_\Phi$ by
the divergence constraint~\eqref{eq:div_constraint}. Squaring gives $\alpha_\Phi^2$.
The integration factor $1/2$ from trajectory averaging along $\SM\to\Mink$ converts
$3\mu^2\alpha_\Phi^2/2\to 3\mu^2\alpha_\Phi^2/4$.
\end{proof}

The ATV prediction for the luminosity--redshift relation was derived in~\cite{STV2017,ATV2026}
by evolving self-similar underdense waves from the end of the radiation epoch to present
time and computing the resulting acceleration. The key observable is the third-order
coefficient $C$ in the standard expansion $H_0 d_L = z + Qz^2 + Cz^3 + O(z^4)$,
which receives contributions from the second-order velocity coefficient $w_2$ in the
STV solution. Under the corrections of Proposition~\ref{prop:EL_reduction}, this
coefficient is modified as described in Theorem~\ref{thm:D} of the next section.

% ============================================================
\section{Redshift Coefficients and Observational Corrections}
\label{sec:redshift}
% ============================================================

\begin{theorem}[Lamphrodyne Correction to the ATV Redshift Coefficient]
\label{thm:D}
The third-order redshift--luminosity coefficient $C$ in the expansion
\begin{equation}
  H_0 d_L = z + Qz^2 + Cz^3 + O(z^4)
  \label{eq:HL_expansion}
\end{equation}
takes the following values in the three frameworks compared in this paper:
\begin{align}
  C_{\Lambda\mathrm{CDM}} &= -0.1804,
  \label{eq:C_LCDM} \\
  C_{\mathrm{ATV}} &= +0.3591,
  \label{eq:C_ATV} \\
  C_{\mathrm{RSVP}}
  &= C_{\mathrm{ATV}}
     + \etaeff
     + \frac{3\mu^2\alpha_\Phi^2}{4}
     + O(\etaeff^2,\mu^3).
  \label{eq:C_RSVP_base}
\end{align}
In the ATV limit $\etaeff,\mu\to 0$ with spherical symmetry enforced,
$C_{\mathrm{RSVP}}\to C_{\mathrm{ATV}}$ exactly. Both correction terms are
positive, so
\begin{equation}
  C_{\mathrm{RSVP}} > C_{\mathrm{ATV}} > 0 > C_{\Lambda\mathrm{CDM}}
  \label{eq:C_ordering}
\end{equation}
whenever $\etaeff > 0$ or $\mu > 0$.
\end{theorem}

\begin{proof}
The ATV derivation~\cite{STV2017,ATV2026} establishes that $C$ is determined at
third order in redshift factor by the value of $w_2$ in the STV-ODE solution
integrated along the trajectory from $\SM$ to $\Mink$. Under
Proposition~\ref{prop:EL_reduction}, the corrections $\delta_\eta\dot{w}_2$
and $\delta_\mu\dot{w}_2$ enter at orders $\etaeff$ and $\mu^2$ respectively.
Integrating along the background trajectory and extracting the $w_2\xi^3$ coefficient
of the velocity field --- which by ATV~\cite{ATV2026} Section~2i determines
the third-order term in the luminosity--redshift expansion --- yields additive
corrections $\etaeff$ and $3\mu^2\alpha_\Phi^2/4$ to $C_{\mathrm{ATV}}$.
The factor $3/4 = 3(\lAn{2}-\lAn{1})^{-1}/2 = 3\cdot(3/2)/2$ arises from the
integration of the perturbation $\delta_\mu\dot{w}_2 = (3\mu^2\alpha_\Phi^2/2)z_2$
along the trajectory. Positivity follows from $\etaeff,\mu,\alpha_\Phi > 0$.
\end{proof}

\begin{remark}[Falsifiability of the base correction]
The null hypothesis $H_0: C_{\mathrm{obs}} = C_{\mathrm{ATV}} = 0.3591$ corresponds
to pure ATV dynamics with $\etaeff = \mu = 0$. RSVP predicts a systematic positive
excess $\Delta C = \etaeff + 3\mu^2\alpha_\Phi^2/4 > 0$. The two corrections are
in principle distinguishable: the $\etaeff$ term is linear in the lamphrodyne
diffusion coefficient and scale-independent, while the $\mu^2$ term depends
quadratically on the scalar--vector coupling and may exhibit scale dependence
through $\alpha_\Phi$. Determining these coefficients observationally requires
redshift--luminosity data to fourth order in $z$, consistent with the ATV
observation~\cite{ATV2026} that the fourth-order coefficient $C_4$ is necessary
to determine whether the Big Bang is self-similar at order $n=2$.
\end{remark}

% ============================================================
\section{The Redshift--Birefringence Cross-Correlation}
\label{sec:bire}
% ============================================================

The central prediction that distinguishes the present framework from both ATV
and minimal $\Lambda$CDM is not birefringence per se --- many theories predict
nonzero polarization rotation --- but the existence of a \emph{nonzero
cross-correlation between anisotropic redshift and anisotropic birefringence},
arising from a common entropy--vorticity source. This cross-correlation,
\begin{equation}
  \Cbz = \bigl\langle\delta\beta(\hat{n})\,\delta z(\hat{n}')\bigr\rangle,
  \label{eq:Cbz_headline}
\end{equation}
is predicted to vanish in both ATV (which has no vorticity sector) and $\Lambda$CDM
(where the birefringence mechanism is decoupled from large-scale structure at first
order), while being generically nonzero in the accessibility relaxation formalism
whenever $\bcurl > 0$. We derive this prediction by relaxing condition
\textup{(RP3)} of the Reduction Principle.

The accessibility formalism generates a natural parity-odd geometric sector.
We derive its minimal form from symmetry principles before postulating a specific
coupling.

\subsection{Derivation of the Parity-Odd Connection Sector}
\label{subsec:parityodd}

Let the full lamphrodyne connection be
$\tilde\Gamma^\mu_{\nu\rho} = \Gamma^\mu_{\nu\rho} + \Delta^\mu_{\nu\rho}$,
where $\Gamma^\mu_{\nu\rho}$ is the Levi-Civita connection of the background
metric~\cite{Wald1984} and $\Delta^\mu_{\nu\rho}$ contains non-Riemannian
corrections induced by gradients and vorticity of $(\Phi,\mathbf{V},S)$.
Decomposing $\Delta$ irreducibly under the action of $\mathrm{SO}(1,3)$ and
spatial inversion yields parity-even and parity-odd sectors.

The parity-odd sector must satisfy three requirements: (i) it is a $(1,2)$-tensor
constructed from $(\Phi,\mathbf{V},S)$ and their first derivatives; (ii) it is
odd under spatial inversion $\mathbf{x}\mapsto-\mathbf{x}$; and (iii) it is linear
in a single scalar quantity at lowest derivative order.

\begin{lemma}[Uniqueness of the minimal parity-odd connection]
\label{lem:parityodd_unique}
Under requirements \textup{(i)--(iii)}, the unique parity-odd $(1,2)$-tensor
linear in a scalar $\Xi$ and its first derivative is
\begin{equation}
  \Delta^{(\mathrm{odd})}_{\mu\nu\rho}
  = \alpha_{\mathrm{ax}}\,\varepsilon_{\mu\nu\rho\sigma}\,\nabla^\sigma\Xi,
  \label{eq:odd_unique}
\end{equation}
where $\varepsilon_{\mu\nu\rho\sigma}$ is the Levi-Civita tensor.
\end{lemma}

\begin{proof}
At first derivative order, a $(1,2)$-tensor odd under spatial inversion must
contain an odd number of spatial indices contracted with $\varepsilon_{\mu\nu\rho\sigma}$.
The unique such structure linear in $\nabla^\sigma\Xi$ is~\eqref{eq:odd_unique}.
All other candidates either vanish by antisymmetry, reduce to~\eqref{eq:odd_unique}
by Bianchi-type identities, or introduce additional derivatives. This is the same
algebraic structure appearing in the Carroll--Field--Jackiw
extension of electrodynamics~\cite{Carroll1990} and in gravitational Chern--Simons
theories~\cite{Wald1984}.
\end{proof}

The lowest-order scalar available in the accessibility formalism that is parity-odd,
nonzero under condition (RP3) relaxed, and couples entropy gradients to vorticity is
\begin{equation}
  \Xi = \nabla^\mu S\,\omega_\mu,
  \qquad
  \omega_\mu = \varepsilon_{\mu\nu\rho\sigma}\,u^\nu\nabla^\rho V^\sigma,
  \label{eq:Xi_def}
\end{equation}
where $\omega_\mu$ is the vorticity of the lamphrodyne flow and $u^\nu$ is the
comoving four-velocity. The quantity $\Xi$ vanishes identically under condition
(RP3) ($\omega = 0$) and is the unique first-order parity-odd scalar with
dimension $[\text{length}]^{-2}$ formed from the available fields. Substituting
into Lemma~\ref{lem:parityodd_unique} gives the connection~\eqref{eq:odd_unique}
and, via the optical transport equation $k^\nu\tilde\nabla_\nu P^\mu = 0$ for the
photon polarization vector~\cite{Wald1984}, the axial projection along null geodesics:

\begin{equation}
  \kA(\eta) = \Pi^\mu_a\Pi^\nu_b\,\epsilon^{ab}\,k^\rho\,\Delta^{(\mathrm{odd})}_{\mu\nu\rho}
  \propto \nabla S\cdot\omega,
  \label{eq:kA_formula}
\end{equation}
and the accumulated CMB polarization rotation $\beta_{\mathrm{obs}} = \int_\gamma\kA\,d\eta$.

\begin{theorem}[Coupled Redshift--Birefringence]
\label{thm:E}
Let $\bcurl > 0$. Then:
\begin{enumerate}[label=\textup{(\alph*)}]
  \item \textbf{Vorticity generation.} The lamphrodyne flow generates vorticity
    $\omega = \nabacc\times\mathbf{V}\neq 0$ generically in the presence of the
    curl coupling $\bcurl|\nabacc\times\mathbf{V}|^2$ in $\Lacc$.
  \item \textbf{Birefringence sourcing.} The parity-odd connection~\eqref{eq:odd_unique}
    produces a nonzero axial projection $\kA \propto \nabla S\cdot\omega\neq 0$,
    generating a polarization rotation $\beta_{\mathrm{obs}}\neq 0$.
  \item \textbf{Redshift correction.} The curl sector contributes an additional term
    \begin{equation}
      \Delta C^{(\bcurl)} = \bcurl\cdot\mathcal{F}(\Phic,S_0) + O(\bcurl^2)
      \label{eq:C_beta_correction}
    \end{equation}
    to $C_{\mathrm{RSVP}}$, where $\mathcal{F}(\Phic,S_0)$ is the linear-response
    functional of $w_2$ to vorticity perturbations along the background trajectory
    (Appendix~\ref{app:F_functional}).
\end{enumerate}
\end{theorem}

\begin{proof}
Part~(a): When $\bcurl > 0$, the term $(\bcurl/2)|\nabacc\times\mathbf{V}|^2$
in $\Lacc$ contributes a non-trivial curl equation in the
$\mathbf{V}$-Euler--Lagrange equation. A non-spherically-symmetric solution will
generically have $\omega\neq 0$; even under \textup{(RP1)}, the nonlinear
interaction of the radial mode with entropy gradients can generate angular momentum
at subleading order in $\xi$.

Part~(b) follows directly from~\eqref{eq:kA_formula}: if $\omega\neq 0$ and $\nabla S\neq 0$
(which holds generically away from perfect isotropy), then $\kA\neq 0$ and the holonomy
$\beta_{\mathrm{obs}} = \int\kA\,d\eta\neq 0$.

Part~(c): The curl term $\bcurl|\nabacc\times\mathbf{V}|^2$ modifies the
$\mathbf{V}$-equation~\eqref{eq:lamp_V} at first order in $\bcurl$. At order $n=2$
of the STV-ODE, the perturbation propagates into $\dot{w}_2$ as $\delta_{\bcurl}\dot{w}_2
= \bcurl\cdot(\partial w_2/\partial\omega)\langle\omega,\nabla S_0\rangle + O(\bcurl^2)$,
which upon integration along the background trajectory defines $\mathcal{F}(\Phic,S_0)$.
\end{proof}

\begin{corollary}[Coupled observables]
\label{cor:E1}
The same parameter $\bcurl > 0$ controls both $\Delta C^{(\bcurl)}\neq 0$
and $\beta_{\mathrm{obs}}\neq 0$. These are not independent observables:
they are controlled by a single coupling constant and satisfy a joint constraint
linking the luminosity--redshift deviation to the polarization rotation magnitude.
\end{corollary}

\begin{corollary}[Joint cross-correlation observable]
\label{cor:E2}
Define the redshift--birefringence cross-correlation
\begin{equation}
  \Cbz = \bigl\langle\delta\beta(\hat{n})\,\delta z(\hat{n}')\bigr\rangle,
  \label{eq:C_bz}
\end{equation}
where $\delta\beta$ is the anisotropic polarization rotation and $\delta z$ the
anisotropic redshift deviation. Then
\begin{equation}
  \bcurl > 0 \implies \Cbz\neq 0,
  \label{eq:C_bz_nonzero}
\end{equation}
through the common entropy--vorticity source $\nabla S\cdot\omega$. No comparable
prediction exists in ATV or minimal $\Lambda$CDM, both of which predict $\Cbz = 0$.
\end{corollary}

\begin{proof}
In ATV, $\bcurl = 0$ enforces $\omega = 0$ and $\kA = 0$ throughout. In $\Lambda$CDM,
the Chern--Simons birefringence mechanism is decoupled from the matter power spectrum
at first order in perturbation theory. In RSVP with $\bcurl > 0$, both $\delta\beta(\hat{n})$
and $\delta z(\hat{n}')$ are sourced by fluctuations of $\nabla S\cdot\omega$ along
lines of sight: $\delta\beta\propto\int\delta(\nabla S\cdot\omega)\,d\eta$ and
$\delta z\propto\int\delta(\nabla_\mu S)\,dx^\mu$. Since both integrals respond to
the same underlying entropy--vorticity field, their cross-correlation is nonzero by the
Cauchy--Schwarz inequality applied to the correlated fluctuations in the RSVP ensemble.
\end{proof}

The proof above uses Cauchy--Schwarz to establish a bound, but a nonzero
cross-correlation requires an additional structural assumption. We make this explicit.

\begin{assumption}[Entropy--vorticity coherence]
\label{ass:coherence}
The stochastic field $X(\hat{n}) = \int_\gamma (\nabla S\cdot\omega)\,d\eta$
possesses a nonvanishing connected two-point function with the entropy gradient:
\begin{equation}
  \bigl\langle\delta X(\hat{n})\,\delta(\nabla S)(\hat{n}')\bigr\rangle_c \neq 0.
  \label{eq:coherence_assumption}
\end{equation}
\end{assumption}

Under Assumption~\ref{ass:coherence}, Corollary~\ref{cor:E2} is strengthened to
an implication rather than a bound. The assumption fails only if entropy gradients
and vorticity are statistically orthogonal at all angular separations, requiring a
special cancellation absent from generic lamphrodyne dynamics.

\subsection{Quantitative Status of the Vorticity Sector}
\label{subsec:vorticity_status}

The parity-odd sector of this section should be understood as a structural
prediction rather than a completed phenomenological model. The existence and joint
structure of the redshift--birefringence coupling $\Cbz\neq 0$ are established
under Assumption~\ref{ass:coherence}, but the \emph{magnitude} of these effects
depends on the response functional $\mathcal{F}(\Phic,S_0)$ of
Appendix~\ref{app:F_functional}. The computation of $\mathcal{F}$ is the principal
quantitative objective of the next stage of this programme; until it is resolved,
the $\omega$-sector contribution to $C_{\mathrm{RSVP}}$ should be treated as
having unknown magnitude relative to the $\etaeff$ and $\mu^2$ corrections.

The complete RSVP prediction for $C$, incorporating all three correction sectors,
is therefore
\begin{equation}
  C_{\mathrm{RSVP}}
  = \underbrace{0.3591}_{C_{\mathrm{ATV}}}
    + \underbrace{\etaeff}_{\kappa_\Phi\text{-sector}}
    + \underbrace{\tfrac{3\mu^2\alpha_\Phi^2}{4}}_{\mu\text{-sector}}
    + \underbrace{\bcurl\cdot\mathcal{F}(\Phic,S_0)}_{\omega\text{-sector}}
    + O(\etaeff^2,\mu^3,\bcurl^2).
  \label{eq:C_full}
\end{equation}
The $\omega$-sector correction is the only one that simultaneously sources birefringence,
and therefore the only one contributing to $\Cbz$.

% ============================================================
\section{Observational Discriminants}
\label{sec:disc}
% ============================================================

Table~\ref{tab:discriminants} summarises the falsifiable predictions distinguishing
RSVP from ATV and $\Lambda$CDM. The three frameworks make identical predictions
for the leading-order Hubble constant $H_0$ and the quadratic coefficient $Q$ in
the luminosity--redshift relation, since these are determined by the order-$n=1$
phase portrait, which is the same in all three. The discriminating power lies
entirely in the third and higher-order coefficients and in correlations between
distinct observables.

\begin{table}[ht]
\centering
\renewcommand{\arraystretch}{1.45}
\begin{tabular}{lccc}
\toprule
Observable & $\Lambda$CDM & ATV & RSVP \\
\midrule
$H_0$, $Q$ (1st, 2nd order) & matched & matched & matched \\
$C$ (3rd-order redshift coeff.) & $-0.1804$ & $+0.3591$ & $>+0.3591$ \\
$\beta_{\mathrm{obs}}$ (birefringence) & via ALP & $0$ & $\neq 0$ (if $\bcurl>0$) \\
$\Cbz$ (redshift--birefringence) & $0$ & $0$ & $\neq 0$ (if $\bcurl>0$) \\
$\dclos$ (closure depth) & --- & $2$ & $2$ \\
$\ell$-dependent depolarization & uniform & absent & present (if $\bcurl>0$) \\
$C_\ell^{\beta\,\mathrm{kSZ}}$ & $0$ & $0$ & $\neq 0$ (if $\bcurl>0$) \\
\bottomrule
\end{tabular}
\caption{Comparison of observational predictions across $\Lambda$CDM, ATV, and
RSVP. The entry ``matched'' indicates that all three frameworks are degenerate
at that order. Entries of $0$ denote predicted null signals. Entries marked
$\neq 0$ denote predicted nonzero signals controlled by the RSVP coupling constant
$\bcurl$. The closure depth $\dclos = 2$ is an algebraic invariant, not directly
observable, but constrains the self-similar structure of the Big Bang.}
\label{tab:discriminants}
\end{table}

The most immediate discriminant between $\Lambda$CDM and both ATV and RSVP is the
sign of $C$: $\Lambda$CDM predicts $C < 0$ while the instability-based frameworks
predict $C > 0$. This sign difference arises because $\Lambda$CDM attributes
acceleration to a constant energy density (dark energy), while ATV and RSVP
attribute it to dynamical instability relaxation, which has different symmetry
properties under time reversal. The sign is in principle measurable from sufficiently
precise supernova Ia surveys; the ATV paper~\cite{ATV2026} notes that the
third-order correction is consistent with an expansion rate that is slowing relative
to dark energy predictions~\cite{DESIBAO2025,DES2025}.

Within RSVP, the discrimination between the pure ATV correction ($C = 0.3591$)
and the full RSVP prediction ($C > 0.3591$) requires fourth-order accuracy in the
luminosity--redshift relation, which lies at the frontier of current observational
precision. The additional discrimination afforded by $\Cbz\neq 0$ is potentially
more accessible: a cross-correlation between anisotropic birefringence maps from
CMB-S4 or LiteBIRD and anisotropic redshift maps from large-scale structure surveys
would constitute a direct probe of $\bcurl$.

\begin{remark}[Current quantitative status of the cross-correlation prediction]
\label{rem:Cbz_status}
The prediction $\Cbz\neq 0$ is a \emph{structural} result: it follows from the
existence of the entropy--vorticity source $\nabla S\cdot\omega$ under
Assumption~\ref{ass:coherence}, and from the common sourcing of both $\delta\beta$
and $\delta z$ by that field. It is not yet a \emph{quantitative} prediction.

The precise amplitude of the $\Cbz$ signal, the angular power spectrum
$C_\ell^{\beta z}$, and its redshift dependence all depend on the response
functional $\mathcal{F}(\Phic,S_0)$ of Appendix~\ref{app:F_functional}, whose
evaluation requires specifying the background entropy profile $S_0(t,r)$ along
the $\SM\to\Mink$ trajectory and computing the linear response of $w_2$ to
vorticity perturbations. Until these computations are complete, the observational
target cannot be specified precisely enough to determine whether CMB-S4 or
LiteBIRD sensitivity is sufficient to detect or constrain the signal.

This gap is explicitly acknowledged as the principal open problem of the programme.
A detection at any significance level would be strong evidence for $\bcurl > 0$;
a null result at the sensitivity levels of those instruments would constrain
$\bcurl\cdot\mathcal{F}$ to be below the noise floor, which would in turn
constrain the vorticity sector of the lamphrodyne functional.
\end{remark}

% ============================================================
\section{Cosmology Beyond Expansion}
\label{sec:beyond}
% ============================================================

The results of this paper collectively support a specific picture of the cosmological
situation. The Friedmann saddle $\SM$ is not the physically realised state of the
universe but rather a metastable configuration --- an accessibility-critical point
of an entropy-relaxation system --- from which the universe has evolved and continues
to evolve. The accelerated expansion that appears anomalous within $\Lambda$CDM is,
on this picture, a generic consequence of relaxation from an unstable configuration,
not a signal of exotic vacuum energy.

This picture has a precise mathematical content. The closure depth $\dclos(\Phic) = 2$
establishes that the full complexity of the instability structure is resolved at the
second perturbative order: the universe's departure from perfect Friedmann behaviour
is a depth-2 phenomenon in the CLIO sense. The RSVP lamphrodyne functional provides
the ontological grounding for what ATV identifies as a dynamical systems structure:
the $\SM\to\Ffam\to\Mink$ phase flow is the cosmological expression of accessibility
relaxation, the trajectory from high-constraint to low-constraint configurations of
the scalar--vector--entropy medium.

Several directions of extension remain open and are noted here explicitly. The first
concerns the full computation of $\mathcal{F}(\Phic,S_0)$, which requires evaluating
the linear response of $w_2$ to vorticity along the background trajectory. This
computation is deferred to a companion paper but is necessary for a quantitative
prediction of the $\omega$-sector contribution to $C_{\mathrm{RSVP}}$ and of the
$\Cbz$ cross-spectrum. The second concerns the radiation-epoch extension: the ATV
analysis is restricted to the matter-dominated ($p=0$) epoch, and a full RSVP
treatment of the radiation epoch ($p = \rho/3$) under the Reduction Principle
would establish whether the instability at the Big Bang itself, rather than
fluctuations from an earlier epoch, is the source of the anomalous acceleration.
The third concerns the $k>0$ Friedmann spacetimes, whose global dynamics are
complicated by the time reversal at maximal expansion and which the ATV paper
defers to subsequent work.

The relationship between RSVP and the ATV framework illustrates a general
methodological principle: rigorous local models within general relativity can serve
as anchor points for a more general field-theoretic ontology, without either reducing
the broader framework to commentary on established physics or requiring the broader
framework to carry full observational burden before its structure is understood.
The ATV analysis provides a precise local shadow of the RSVP lamphrodyne relaxation
geometry; the RSVP framework provides that shadow with an ontological substrate
and a programme of corrections.

The deepest open question concerns what lies beyond the ATV limit in the RSVP
cosmology. If the universe is genuinely a lamphrodyne relaxation system rather than
a Friedmann spacetime with corrections, then the full $(\Phi,\mathbf{V},S)$
dynamics govern structure formation, void growth, CMB anisotropy, and birefringence
in an integrated manner. The observational discriminants of Table~\ref{tab:discriminants}
are the first predictions from this regime that are accessible to current or
near-future experiments. A detection of $\Cbz\neq 0$ would constitute the
strongest available evidence that the parity-odd sector of the universe's geometry
is sourced by entropy--vorticity coupling of the kind the RSVP framework predicts,
rather than by a propagating pseudoscalar field of the kind minimal axion models predict.

\subsection{Outstanding Problems and the Scope of the Current Framework}
\label{subsec:open_problems}

For clarity and intellectual honesty, we enumerate the principal open problems
that must be resolved before the accessibility relaxation formalism can be
considered a mature cosmological theory. These open problems do not affect the
internal consistency of the ATV-sector results in
Sections~\ref{sec:friedmann}--\ref{sec:closure}, but they do define the next
stage of development.

The most urgent technical problem is the computation of
$\mathcal{F}(\Phic,S_0)$, which is required to convert the structural prediction
$\Cbz\neq 0$ into a quantitative observational target. Until $\mathcal{F}$ is
evaluated, the vorticity sector contributes a term of unknown magnitude to
$C_{\mathrm{RSVP}}$.

The extension to the radiation-dominated epoch ($p = \rho/3$) is necessary to
establish whether the RSVP instability survives in the high-pressure environment
of the early universe and whether the radiation-epoch dynamics alter the
closure-depth invariant or the eigenvalue spectrum.

The derivation of the coupling constants $(\kappa, \alpha, \bcurl, \mu)$ from
a microscopic theory of accessibility, rather than from symmetry counting and
phenomenological consistency, would convert the effective-theory description of
Section~\ref{rem:eft_status} into a fundamental derivation and would allow
precise predictions without free parameters.

A complete nonlinear stability analysis of the full scalar--vector--entropy system
beyond the linearised Hessian regime is needed to verify that the depth-two
closure theorem survives into the nonlinear dynamical regime, and that no
nonlinear mode coupling introduces additional stable directions at orders $n\geq 3$.

Finally, the gauge invariance of $\dclos$ established in
Proposition~\ref{prop:closure_invariance} should be extended to cover
diffeomorphism equivalences of the full lamphrodyne system, not only of
the truncation hierarchy, to confirm that the invariant is a geometric
property of the accessibility bundle rather than an artefact of the
self-similar coordinate system.

% ============================================================
\appendix
% ============================================================

\section{The Functional $\mathcal{F}(\Phi_c,S_0)$}
\label{app:F_functional}

The functional $\mathcal{F}(\Phic,S_0)$ appearing in Theorem~\ref{thm:E}(c) and
equation~\eqref{eq:C_full} encodes the linear response of the second-order velocity
coefficient $w_2$ to vorticity perturbations along the background lamphrodyne
trajectory from $\SM$ to $\Mink$. It is defined by
\begin{equation}
  \mathcal{F}(\Phic,S_0)
  = \int_{\SM}^{\Mink}
    \bigl\langle\omega_0(t),\,\nabla S_0\bigr\rangle\,
    \frac{\partial w_2}{\partial\omega}\,
    dt,
  \label{eq:F_functional_def}
\end{equation}
where $\omega_0$ is the background vorticity, $S_0$ the background entropy profile,
and $\partial w_2/\partial\omega$ is the first-order variation of $w_2$ with respect
to a vorticity perturbation $\delta\omega$ at fixed $(z_2,w_0)$.

To compute $\partial w_2/\partial\omega$ at leading order, one perturbs the ATV
STV-ODE at order $n=2$ by the additional source term
$\delta_{\bcurl}\dot{w}_2 = \bcurl\langle\delta\omega,\nabla S_0\rangle$ and
integrates along the unperturbed trajectory. The result depends on the background
entropy gradient profile $\nabla S_0$ evaluated along $\SM\to\Mink$, which
in turn depends on the initial conditions for the entropy field $S_0$ at the
end of the radiation epoch.

A full computation of $\mathcal{F}$ requires: specification of the background entropy
profile $S_0(t,r)$ consistent with RSVP dynamics over the matter-dominated epoch;
evaluation of the Green's function for the perturbed $w_2$ equation; and integration
against $\langle\omega_0,\nabla S_0\rangle$ along the trajectory. This computation
is deferred to a companion paper. The sign of $\mathcal{F}$ is expected to be positive,
consistent with the RSVP principle that entropy gradients amplify accessible trajectory
volume, which would further increase $C_{\mathrm{RSVP}}$ above $C_{\mathrm{ATV}}$.

\section{Hierarchy of Correction Terms}
\label{app:hierarchy}

The three correction sectors of equation~\eqref{eq:C_full} are associated with
distinct physical mechanisms and have distinct observational signatures:

The $\kappa_\Phi$-sector correction $\etaeff = (2/3)\kappa_\Phi$ is scale-independent
and does not source birefringence. It is the lamphrodyne smoothing of $\Phi$ in the
RSVP Master System: at any nonzero diffusion coefficient $\kappa_\Phi > 0$, the scalar
field relaxes toward spatial uniformity, which modifies the effective pressure on the
trajectory bundle and shifts $w_2$ by a universal additive constant.

The $\mu$-sector correction $3\mu^2\alpha_\Phi^2/4$ is quadratic in the scalar--vector
coupling and arises from integrating out the longitudinal flow mode via the divergence
constraint~\eqref{eq:div_constraint}. It is generically present whenever $\Phi$ and
$\mathbf{V}$ are coupled and does not source birefringence.

The $\omega$-sector correction $\bcurl\cdot\mathcal{F}(\Phic,S_0)$ is the only one
that simultaneously sources birefringence and contributes to $\Cbz$. It is absent
in any spherically symmetric reduction and vanishes in the ATV limit, confirming
that the ATV framework is a consistent truncation of RSVP in which torsion and
vorticity are set to zero from the outset.

\section{ATV Phase Portrait in RSVP Coordinates}
\label{app:phase_portrait}

The ATV phase portrait at order $n=1$, shown in Figure~1 of~\cite{ATV2026},
describes the flow of solutions from the source $U=(0,0)$ through the saddle $\SM$
to the attractor $\Mink$. In RSVP coordinates, this portrait is the projection of
the full lamphrodyne flow onto the $(z_2,w_0)$ plane under Reduction
Principle~\ref{prin:reduction}. The three rest points correspond to:
\begin{itemize}[leftmargin=1.5em]
  \item $U = (0,0)$: the empty configuration with vanishing scalar field and zero
    flow velocity. Eliminated by the time-since-the-Big-Bang gauge fixing, it
    corresponds in RSVP to the zero-accessibility vacuum, which is unphysical as
    a cosmological initial condition.
  \item $\SM = (4/3, 2/3)$: the accessibility-critical configuration of
    Theorem~\ref{thm:B}. In RSVP, this is a metastable compressed state of the
    scalar--vector medium: high $\Phi$, ordered flow, and near-uniform entropy.
    Its instability reflects the generic tendency of accessibility-gradient systems
    to relax away from constrained configurations.
  \item $\Mink = (0,1)$: the low-curvature asymptotic accessibility basin.
    In RSVP, this corresponds to maximal accessibility: vanishing scalar gradient,
    isotropic flow, and entropy equilibrium. It is the attractor of the lamphrodyne
    relaxation and represents not expansion into emptiness but the gradual flattening
    of constraint structure within the medium.
\end{itemize}
The underdense trajectory connecting $\SM$ to $\Mink$ is the family $\Ffam$ of
ATV underdense solutions, which in RSVP corresponds to the one-parameter family of
entropy-release trajectories admissible from the compressed initial configuration.
Generic members of $\Ffam$ accelerate away from Friedmann spacetimes at intermediate
times and return to asymptotic Friedmann behaviour at late times --- a pattern
that RSVP attributes to the nonlinear interaction of accessibility-gradient relaxation
with the constraint curvature of the medium at second perturbative order.

% ============================================================
% REFERENCES
% ============================================================

\begin{thebibliography}{99}

\bibitem{Smoller2012}
J.~Smoller and B.~Temple,
``General relativistic self-similar waves that induce an anomalous acceleration
into the standard model of cosmology,''
\textit{Mem.\ Amer.\ Math.\ Soc.} \textbf{218} (2012).

\bibitem{STV2017}
J.~Smoller, B.~Temple, and Z.~Vogler,
``An instability of the standard model of cosmology creates the anomalous
acceleration without dark energy,''
\textit{Proc.\ R.\ Soc.\ A} \textbf{473}, 20160887 (2017).

\bibitem{ATV2026}
C.~Alexander, B.~Temple, and Z.~Vogler,
``The instability of critical and underdense Friedmann spacetimes at the Big Bang
as an alternative to dark energy,''
\textit{Proc.\ R.\ Soc.\ A} \textbf{482}, 20250912 (2026).
\url{https://doi.org/10.1098/rspa.2025.0912}

\bibitem{Jacobson1995}
T.~Jacobson,
``Thermodynamics of spacetime: the Einstein equation of state,''
\textit{Phys.\ Rev.\ Lett.} \textbf{75}, 1260--1263 (1995).

\bibitem{Verlinde2011}
E.~Verlinde,
``On the origin of gravity and the laws of Newton,''
\textit{JHEP} \textbf{2011}(04), 029 (2011).

\bibitem{Padmanabhan2010}
T.~Padmanabhan,
``Thermodynamical aspects of gravity: new insights,''
\textit{Rep.\ Prog.\ Phys.} \textbf{73}, 046901 (2010).

\bibitem{HawkingEllis1973}
S.~W.~Hawking and G.~F.~R.~Ellis,
\textit{The Large Scale Structure of Space-Time},
Cambridge University Press (1973).

\bibitem{Wald1984}
R.~M.~Wald,
\textit{General Relativity},
University of Chicago Press (1984).

\bibitem{MisnerThorneWheeler1973}
C.~W.~Misner, K.~S.~Thorne, and J.~A.~Wheeler,
\textit{Gravitation},
W.~H.~Freeman (1973).

\bibitem{Weinberg1972}
S.~Weinberg,
\textit{Gravitation and Cosmology},
Wiley, New York (1972).

\bibitem{EllisMaartensMacCallum2012}
G.~F.~R.~Ellis, R.~Maartens, and M.~A.~H.~MacCallum,
\textit{Relativistic Cosmology},
Cambridge University Press (2012).

\bibitem{Peebles1993}
P.~J.~E.~Peebles,
\textit{Principles of Physical Cosmology},
Princeton University Press (1993).

\bibitem{Arnold1989}
V.~I.~Arnold,
\textit{Mathematical Methods of Classical Mechanics},
2nd ed., Springer, New York (1989).

\bibitem{Hirsch1974}
M.~W.~Hirsch and S.~Smale,
\textit{Differential Equations, Dynamical Systems, and Linear Algebra},
Academic Press, New York (1974).

\bibitem{Pazy1983}
A.~Pazy,
\textit{Semigroups of Linear Operators and Applications to Partial Differential
Equations}, Springer, New York (1983).

\bibitem{Lions1969}
J.-L.~Lions,
\textit{Quelques m\'{e}thodes de r\'{e}solution des probl\`{e}mes aux limites non
lin\'{e}aires}, Dunod, Paris (1969).

\bibitem{Carroll1990}
S.~M.~Carroll, G.~B.~Field, and R.~Jackiw,
``Limits on a Lorentz- and parity-violating modification of electrodynamics,''
\textit{Phys.\ Rev.\ D} \textbf{41}, 1231 (1990).

\bibitem{Minami2020}
Y.~Minami and E.~Komatsu,
``New extraction of cosmic birefringence from Planck 2018 polarization data,''
\textit{Phys.\ Rev.\ Lett.} \textbf{125}, 221301 (2020).

\bibitem{CMBS4}
CMB-S4 Collaboration,
``CMB-S4 Science Book,''
arXiv:1610.02743 [astro-ph.CO] (2016).

\bibitem{LiteBIRD}
LiteBIRD Collaboration,
``Probing cosmic inflation with the LiteBIRD cosmic microwave background
polarization survey,''
\textit{PTEP} \textbf{2023}, 042F01 (2023).

\bibitem{DESIBAO2025}
K.~Lodha et al.,
``Extended dark energy analysis using DESI DR2 BAO measurements,''
\textit{Phys.\ Rev.\ D} \textbf{112}, 083511 (2025).

\bibitem{DES2025}
T.~M.~C.~Abbott et al.\ (DES Collaboration),
``Dark energy survey: implications for cosmological expansion models from the
final DES baryon acoustic oscillation and supernova data,''
\textit{Phys.\ Rev.\ D} \textbf{113}, 063530 (2025).

\end{thebibliography}

\end{document}

